\section*{अध्याय \ref{ch:b3:gravity-geodesy} --- गुरुत्व और पृथ्वी की आकृति}

\begin{solution}{exo:b3:gravity-geodesy:1}
$b = a(1 - f) = 6378137\times0.99664719 = \qty{6356752.3}{m}$ और
$a - b = \qty{21384.7}{m}$, लगभग \qty{21.4}{km}। क्लेरो:
$f = (3J_2 + m)/2 = (0.00324789 + 0.00344979)/2 = \num{0.0033488}$ या
$1/298.61$, जबकि $1/298.26$ के मुकाबले \qty{0.12}{\percent} छोटा है, यह अंतर उन
दूसरी-कोटि के पदों का आकार है, जिन्हें यह संबंध उपेक्षित करता है।
\end{solution}

\begin{solution}{exo:b3:gravity-geodesy:2}
$\sigma = \qty{6.0e13}{kg}/\qty{5.0e11}{m^2} = \qty{120}{kg.m^{-2}}$, अतः
$h_w = 120/1000 = \qty{0.12}{m}$: यानी \qty{12}{cm} जल। महासागर में पहुँचाए जाने पर
$60/360 = \qty{0.17}{mm}$ समुद्र-तल में परिवर्तन।
\end{solution}

\begin{solution}{exo:b3:gravity-geodesy:3}
भूमध्य रेखा पर $\gamma = \qty{9.780327}{m.s^{-2}}$। \ang{45} पर,
$\sin^2\varphi = 0.5$ और $\sin^22\varphi = 1$: $\gamma = 9.780327\times1.0026454
= \qty{9.806200}{m.s^{-2}}$। ध्रुवों पर, $9.780327\times1.0053024 =
\qty{9.832186}{m.s^{-2}}$। अंतर \qty{0.051859}{m.s^{-2}} है, या
\qty{5186}{\milli\gal}.
\end{solution}

\begin{solution}{exo:b3:gravity-geodesy:4}
$D = ET_e^3/(12(1 - \nu^2)) = 10^{11}T_e^3/11.25$: \qty{8.89e21}{N.m} का मान,
\qty{10}{km} के लिए, \qty{2.40e23}{N.m} का मान \qty{30}{km} के लिए। जहाँ $(\rho_m - \rho_w)g =
\qty{22269}{N.m^{-3}}$, $\alpha = (4D/22269)^{1/4}$: \qty{35.5}{km} और
\qty{81.0}{km}। खाई $3\pi\alpha/4$ पर समाप्त होती है, \qty{84}{km} और \qty{191}{km} पर;
\omterm{def:b3:gravity-geodesy:alpha}{अग्र-उभार} के शिखर $\pi\alpha$ पर हैं, \qty{112}{km} और \qty{255}{km} पर।
\end{solution}

\begin{solution}{exo:b3:gravity-geodesy:5}
(क) $T = -\gamma N = -9.80\times50 = \qty{-490}{J.kg^{-1}}$। (ख) भूआभ
दीर्घवृत्तज के सापेक्ष $10/100\,000 = \num{1.0e-4}$ रेडियन झुका है, और ओलंब रेखा,
जो भूआभ के लंबवत है, दीर्घवृत्तज के अभिलंब के सापेक्ष उसी कोण से झुकी है:
$\num{1.0e-4}\times206\,265 = 20.6''$.
\end{solution}

\begin{solution}{exo:b3:gravity-geodesy:6}
$H = J_2/(C/Ma^2) = \num{1.08263e-3}/0.3308 = \num{0.0032728}$ और
$1/H = 305.55$। परिकल्पना से $J_2$ समान है; एक समान पृथ्वी में
$C/(Ma^2) = 0.4$ होता, और इसलिए $H = J_2/0.4 = \num{0.0027066}$: पुरस्सरण
धीमा होता। $J_2$ और $H$ दोनों को मापना ही वह बात प्रकट करता है कि द्रव्यमान
केंद्र की ओर संकेंद्रित है।
\end{solution}

\begin{solution}{exo:b3:gravity-geodesy:7}
$r = \qty{6878.137}{km}$; $n = \sqrt{GM/r^3} = \qty{1.10678e-3}{rad.s^{-1}}$, जिसका
आवर्तकाल $2\pi/n = \qty{94.6}{min}$ है। $\dot\Omega = -1.5\times\num{1.10678e-3}
\times\num{1.08263e-3}\times(6378.137/6878.137)^2\times\cos\ang{89} =
\qty{-2.70e-8}{rad.s^{-1}}$, या \ang{-0.134} प्रतिदिन, पश्चिम की ओर। एक लगभग-ध्रुवीय कक्षा
हर अक्षांश के ऊपर से गुज़रती है, इसलिए पूरी पृथ्वी का मानचित्रण होता है; यह भी
धीरे-धीरे पुरस्सरण करती है, ताकि भू-पथ नियमित रूप से भर जाएँ।
\end{solution}

\begin{solution}{exo:b3:gravity-geodesy:8}
$e^{-2\pi\times500/100} = e^{-31.4} = \num{2.3e-14}$: \qty{2e-13}{\milli\gal},
कुछ नहीं। $e^{-2\pi\times500/1000} = e^{-3.14} = 0.0432$:
\qty{0.43}{\milli\gal}.
\end{solution}

\begin{solution}{exo:b3:gravity-geodesy:9}
$3\pi\alpha/4 = \qty{120}{km}$ लेने पर $\alpha = \qty{50.9}{km}$ प्राप्त होता है।
$D = (\rho_m - \rho_w)g\alpha^4/4 = 22269\times(50\,930)^4/4 =
\qty{3.75e22}{N.m}$, और $T_e = (11.25D/E)^{1/3} = \qty{16.2}{km}$। प्लेट
खाई से $\pi\alpha/2 = \qty{80}{km}$ पर अपने आनमन-रहित स्तर को पार करती है।
\end{solution}

\begin{solution}{exo:b3:gravity-geodesy:10}
(क) $u = x/\alpha$ लेने पर, $\mathrm d^4/\mathrm dx^4 = \alpha^{-4}\,\mathrm
d^4/\mathrm du^4$। $\cos u + \sin u = \sqrt2\cos(u - \pi/4)$ लिखने पर
फलन $\sqrt2w_0\,\mathrm{Re}\,e^{(-1 + i)u - i\pi/4}$ है, और
$(-1 + i)^4 = -4$: अतः $w'''' = -4w/\alpha^4$। तब $Dw'''' + (\rho_m -
\rho_i)gw = (-4D/\alpha^4 + (\rho_m - \rho_i)g)w = 0$, क्योंकि $\alpha^4 =
4D/((\rho_m - \rho_i)g)$। (ख) $w'' = -(2w_0/\alpha^2)e^{-u}(\cos u - \sin u)$, अतः
$M = -Dw'' = (2Dw_0/\alpha^2)e^{-u}(\cos u - \sin u)$: परिमाण में सबसे बड़ा
$u = 0$ पर है, $\tan u = 1$ पर शून्य, $x = \pi\alpha/4$ पर। (ग) भार के नीचे प्लेट
ऊपर की ओर अवतल है, इसकी ऊपरी सतह संपीड़न में; $\pi\alpha/4$ से परे
यह ऊपर की ओर उत्तल है, इसकी ऊपरी सतह खिंची हुई। वहाँ $M' \propto
-e^{-u}\cos u$ का मान $u = \pi/2$ पर शून्य होता है: खिंचाव सबसे अधिक
$x = \pi\alpha/2$ पर है, खाई की \omterm{def:b3:gravity-geodesy:alpha}{अग्र-उभार} की ओर वाली बाहरी ढलान पर, जहाँ
$|M|$ भार के नीचे अपने मान का $e^{-\pi/2} = \qty{21}{\percent}$ है। यही वह स्थान है,
जहाँ प्लेट के शीर्ष पर सामान्य भ्रंश और दरारें खुलेंगे।
\end{solution}

\begin{solution}{exo:b3:gravity-geodesy:11}
$\sqrt{1 + \eta} = 2.5(1 - 1.5\times0.3308) = 1.2595$, अतः $1 + \eta = 1.5863$
और $\tfrac52m/f = 2.5863$: $f_h = 2.5\times0.00344979/2.5863 = \num{0.0033346}$,
$1/f_h = 299.9$। अंतर $f - f_h = \num{1.82e-5}$ है; $a$ से गुणा करने पर,
भूमध्य रेखा द्रवस्थैतिक आकृति से \qty{116}{m} ऊपर है। पृथ्वी
पूर्णतः साम्यावस्था में द्रव नहीं है: इसका आंतरिक भाग प्रतिबल सहता है, जो स्थलमंडल की सामर्थ्य से
या प्रावार के प्रवाह और घनत्व अंतरों से उत्पन्न होता है, जो
आकृति को उससे थोड़ा दूर रखता है।
\end{solution}

\begin{solution}{exo:b3:gravity-geodesy:12}
कक्षीय पात-बिंदु को $2\pi/(365.242\times86\,400) =
\qty{1.991e-7}{rad.s^{-1}}$ से आगे बढ़ना होगा। $r = \qty{7078.137}{km}$ पर, $n =
\qty{1.0602e-3}{rad.s^{-1}}$ और $\tfrac32nJ_2(a/r)^2 = \num{1.398e-6}$, अतः
$\cos i = -\num{1.991e-7}/\num{1.398e-6} = -0.1424$ और $i = \ang{98.2}$: एक
थोड़ी प्रतिगामी, सूर्य-तुल्यकालिक कक्षा।
\end{solution}

\begin{solution}{pb:b3:gravity-geodesy:1}
\textbf{1.} $m = (\num{7.292115e-5})^2\times6378137^2\times6356752.3/
\num{3.986005e14} = \num{0.00344979}$.

\textbf{2.} $f = (3J_2 + m)/2 = \num{0.0033488}$, $1/f = 298.61$.

\textbf{3.} $\sqrt{1 + \eta} = 2.5(1 - 1.5\times0.3308) = 1.2595$,
$\tfrac52m/f_h = 2 + 0.5863$, अतः $f_h = \num{0.0033346}$ और $1/f_h = 299.88$।

\textbf{4.} क्लेरो मान, $1/298.61$, प्रेक्षित $1/298.26$ से
\qty{0.12}{\percent} के भीतर है; जलस्थैतिक मान, $1/299.88$, इससे
$\Delta f = \num{1.82e-5}$ के अंतर से छोटा है, और $a\,\Delta f = \qty{116}{m}$: विषुवत रेखा
जलस्थैतिक आकृति से \qty{116}{m} ऊपर स्थित है।

\textbf{5.} $H = \num{1.08263e-3}/0.3308 = \num{0.0032728}$ ($1/305.6$).

\textbf{6.} $\tfrac12\times\qty{80}{km}\times\qty{6}{km} = \qty{240}{km^2} =
\qty{2.40e8}{m^2}$.

\textbf{7.} $V_0 = (2800 - 1030)\times9.81\times\num{2.40e8} =
\qty{4.17e12}{N.m^{-1}}$.

\textbf{8.} एरी: $(\rho_m - \rho_v)r = (\rho_v - \rho_w)h$, अतः $r =
6\times1770/500 = \qty{21.2}{km}$ शिखा के नीचे।

\textbf{9.} समुद्री पेंदी कटक के आधार से काफ़ी दूर तक विक्षेपित होती है: भार
स्तंभ-दर-स्तंभ प्रतिपूरित नहीं होता, बल्कि चार गुना चौड़े क्षेत्र पर प्लेट के
आनमन से प्रतिपूरित होता है। प्लेट में सामर्थ्य है: \omterm{def:b3:gravity-geodesy:regional}{क्षेत्रीय समस्थितिक प्रतिपूर्ति}।

\textbf{10.} $\alpha = 165/(3\pi/4) = \qty{70.0}{km}$.

\textbf{11.} $\alpha = 220/\pi = \qty{70.0}{km}$: वही, जैसा एक एकल प्रत्यास्थ
प्लेट के लिए अपेक्षित है।

\textbf{12.} $D = (\rho_m - \rho_w)g\alpha^4/4 = 2270\times9.81\times(\num{7.00e4})^4/4
= \qty{1.34e23}{N.m}$.

\textbf{13.} $T_e = (12(1 - \nu^2)D/E)^{1/3} = (11.25\times\num{1.34e23}/
\num{1e11})^{1/3} = \qty{24.7}{km}$.

\textbf{14.} $w_0 = V_0\alpha^3/(8D) = \num{4.17e12}\times(\num{7.00e4})^3/
(8\times\num{1.34e23}) = \qty{1334}{m}$.

\textbf{15.} $w_0e^{-\pi} = 1334\times0.0432 = \qty{58}{m}$.

\textbf{16.} $2\times1334\times\num{7.00e4}\times2270\times9.81 =
\qty{4.16e12}{N.m^{-1}}$, $V_0$ के बराबर है, $D$ के पूर्णांकन की सीमा के भीतर।

\textbf{17.} \qty{24.7}{km} हवाई कटक की परास के भीतर आता है। प्लेट
अक्ष के नीचे \qty{1.3}{km} धँसती है, बजाय इसके कि \qty{21}{km} का मूल बने: वही
उत्प्लावकता \qty{330}{km} लंबी परिखा में फैली रहती है।

\textbf{18.} $\lambda = 2\pi\times\qty{70.0}{km} = \qty{440}{km}$,
$k = 1/\alpha = \qty{1.43e-5}{m^{-1}}$.

\textbf{19.} $\sigma = (2900 - 1030)\times500 = \qty{9.35e5}{kg.m^{-2}}$, अतः
$2\pi G\sigma = \qty{3.92e-4}{m.s^{-2}} = \qty{39.2}{\milli\gal}$ समुद्री पेंदी पर;
$e^{-k\times\qty{5}{km}} = 0.931$ से गुणा करने पर, समुद्र-तल पर \qty{36.5}{\milli\gal}
(जहाँ प्लेट नीचे दबाई गई है वहाँ ऋणात्मक)।

\textbf{20.} $\sigma = (3300 - 2900)\times500 = \qty{2.0e5}{kg.m^{-2}}$:
\qty{8.39}{\milli\gal}, $e^{-k\times\qty{12}{km}} = 0.843$ से गुणा करने पर,
\qty{7.07}{\milli\gal}। कुल योग \qty{43.6}{\milli\gal} है।

\textbf{21.} $e^{-k\times\qty{400}{km}} = e^{-5.71} = 0.0033$: \qty{0.14}{\milli\gal}.

\textbf{22.} $e^{-2\pi\times400/80} = e^{-31.4} = \num{2.3e-14}$: उपग्रह
कटक का अपना कुछ भी नहीं देखता, केवल विस्तृत आनमन देखता है, और वह भी
तीन-सौ गुना क्षीण होकर; कटक की अपनी गुरुत्व विसंगति समुद्र में मापी जानी चाहिए।

\textbf{23.} $n \approx 2\pi\times6371/440 = 91$.

\textbf{24.} प्रश्न 10 से 13 के आधार पर, ओरु कटक के नीचे की प्लेट की
\omterm{def:b3:gravity-geodesy:flexure}{प्रभावी प्रत्यास्थ मोटाई} है
\textbf{\qty[text-series-to-math=true]{24.7}{km}}।
\end{solution}
